\chapter{The Tube: BBC BASIC}\label{cha:tube}

The BBC Micro's most elegant idea was the Tube: a fast port through which a
\emph{second processor} --- another CPU with its own memory --- could take
over the computation while the Beeb kept the keyboard, the screen and the
discs. The K4510 has a Tube of its own at \texttt{\$D800}, and the first
thing fitted to it is Richard Russell's BBC BASIC, running on the host
machine with a flat 256\,MB of its own.

\begin{type}
BBC
\end{type}

You get the \verb|>| prompt of a BASIC with real power behind it, and it
is \emph{fast} --- the co-processor is not cycle-matched to 1985:

\begin{type}
HIMEM=PAGE+250*1024*1024
DIM space 200*1024*1024
\end{type}

both just work. Type \verb|*QUIT| to hand the console back to the shell.

\section{Graphics and sound}
The console cannot show what it cannot see, so the Tube speaks up: MODE,
GCOL, PLOT, MOVE, DRAW, CIRCLE and the palette arrive at the machine as
escape sequences, and the Tube's gatekeeper (the \emph{Tube ULA}) executes
them on the VICKY blitter. \verb|SOUND| and \verb|ENVELOPE|-less music go
the same way, onto the machine's four-channel sound sequencer and out
through the OPL2 --- the classic Beeb \texttt{SOUND chan,amp,pitch,dur},
queued per channel like the original. Channel~0, which was the Beeb's
noise, is a feedback-heavy FM patch; 1 to 3 are plain two-operator tones.
The patch is written again on every note, so a program that has zeroed
the chip for its own purposes does not silence the sequencer.

\begin{type}
CD /LANG/BBCBASIC
BBC
LOAD "EX/KALEID.BBC"
RUN
\end{type}

\screen{shots/bbc.png}{\texttt{KALEID.BBC}: MODE 2 kaleidoscope, drawn by
BBC BASIC on the Tube, painted by VICKY.}

The \pth{/LANG/BBCBASIC/EX} directory carries a period programme
selection:
\texttt{KALEID}, \texttt{CIRCLES}, \texttt{ROSES}, \texttt{MOUNTAIN},
\texttt{BOUNCE}, \texttt{CLOCK} (a live analogue clock face) and
\texttt{TUNE} --- Fr\`ere Jacques as a three-voice round, the sound
sequencer's party piece.

\section{Sprites}
VICKY has 128 hardware sprites, and BBC BASIC reaches them the way RISC OS
reached its own: \verb|VDU 23,27| selects and shapes, \verb|PLOT &ED|
places. There is no sprite file. A sprite is \emph{captured} from the
bitmap after BBC BASIC has drawn it with the words it already knows:

\begin{type}
MODE 2
GCOL 0,1:CIRCLE FILL 40,40,30
MOVE 8,6:VDU 23,27,1,0,32,32|
CLG
PLOT &ED,600,500
\end{type}

\small
\begin{tabular}{@{}l >{\raggedright\arraybackslash}p{0.54\linewidth}@{}}
\verb|VDU 23,27,0,n|| & select sprite $n$ (0--127) for PLOT \\
\verb|VDU 23,27,1,n,w,h|| & capture $n$, $w\times h$ pixels (8/16/32/64), bottom-left at the graphics cursor \\
\verb|VDU 23,27,2,n|| & hide $n$ \\
\verb|VDU 23,27,3,n,f|| & flip: bit 0 horizontal, bit 1 vertical \\
\verb|VDU 23,27,4,n,z|| & depth: drawn after layer $z$ (default 1, over the bitmap) \\
\verb|VDU 23,27,5,n,m|| & sprite $n$ shows sprite $m$'s picture \\
\verb|PLOT &ED,x,y| & show the selected sprite, bottom-left at $x,y$ \\
\end{tabular}
\normalsize

A placed sprite is a register write: moving sixty-four of them every frame
costs the interpreter nothing but the PLOTs. \texttt{SPRITES.BBC},
\texttt{INVADERS.BBC} and \texttt{TUNNEL.BBC} beside them are
the demonstrations (the last one moves nothing at all --- it cycles the
palette with \verb|VDU 19|).

\section{The terminal: JIM}
The console you see when the Tube is running is not the ROM's own terminal
but \emph{JIM}, at \texttt{\$DA00} --- the Beeb's third I/O page, given a
job here. JIM is a VT100 with the ANSI colours and the VT220 editing
sequences (insert and delete character and line, erase character, scroll
regions, origin mode, DEC line drawing, cursor-position and identity
reports), built into the machine the way a real 8-bit computer got a
serious terminal: as a card, not a program. It draws on the same text
screen as the ROM console, inside the window the ROM gives it, so the two
share one screen and one cursor. Anything that needs a terminal writes its
byte stream to \texttt{\$DA00} and reads its answers from \texttt{\$DA02};
keys pushed through \texttt{\$DA03} come back as the VT sequences the far
end expects (arrows, Home/End, PgUp/PgDn, F1--F12, Delete as \texttt{\$7F}).
The register map is in \texttt{core/term.h}.

\screen{shots/turbo.png}{Turbo Pascal~3 (\texttt{H:} user~3) on JIM: the
menu bar in reverse video, the compiler's own screen handling on a VT100
that is a chip.}

For CP/M this settles the question every program asks at install time:
tell WordStar's \texttt{WSCHANGE}, Turbo Pascal's \texttt{TINST}, ZDE and
the rest that the terminal is a \emph{VT100} (or \emph{ANSI}: the same
sequences). Turbo Pascal~3 on \texttt{H:} user~3 comes up with its menu
already right; WordStar~4 on \texttt{E:} wants one pass of
\texttt{WSCHANGE}. BBC BASIC's console edition speaks the same language
natively, so \verb|COLOUR|, \verb|CLS| and \verb|PRINT TAB(x,y)| land
where they should. Bytes \texttt{\$80}--\texttt{\$FF} are drawn as CP437
glyphs, which is what BBS ANSI art is made of.

\section{Star commands}
A line starting with \verb|*| goes to the machine: \verb|*DIR|, \verb|*CD|,
any shell command --- and BBC BASIC's own file words (\verb|LOAD|,
\verb|SAVE|) read and write the machine's filesystem directly, because the
co-processor lives inside \texttt{fs/} too. The Tube starts in the shell's
current directory, so \verb|CD /LANG/BBCBASIC| then \verb|BBC| lets
\verb|LOAD "EX/KALEID.BBC"| work without spelling the whole path; the machine's
filesystem is the co-processor's whole world --- \texttt{fs/} is shown as
\texttt{/}, the host tree above it hidden, and \verb|*CD ..| stops at the
root.

\begin{aside}
\textbf{Edges, honestly:} \texttt{POINT(} and \texttt{TINT} are not
implemented; \texttt{GCOL} modes 1--4 draw plain; \texttt{VDU 5} text
prints as text.
\end{aside}
